% !TEX root = ../main.tex
\chapter{Differentiable structures: definition and classification}\label{ch:l07}
\begin{paracol}{2}

\begin{sources}
\begin{tabularx}{\linewidth}{@{}l l L@{}}
\toprule
\textbf{Video} & \textbf{Length} & \textbf{Content}\\
\midrule
\seg{7}{1}{Lecture 7: Differentiable structures: definition and classification} & 1:14:34 & Refining the maximal $C^0$ atlas: $C^k$, smooth, analytic and complex atlases;
the theorem that a maximal $C^k$ atlas ($k\ge1$) contains a unique smooth one; differentiable manifolds; two incompatible smooth atlases of the real line; differentiability of maps between manifolds and its
well-definedness; diffeomorphisms; levels of structure on a set (and student questions); how many smooth structures a topological manifold carries (dimensions 1--3, above 4, exactly 4);
tangent spaces without embedding; outlook to tensor spaces and modules.\\
\bottomrule
\end{tabularx}

\smallskip
\textbf{Transcript} \file{materials/transcripts/C07-V01.txt} (YouTube automatic captions).
\end{sources}

\begin{reading}[Where this lecture sits in the other notes]
\emph{Gravity and Light}, Lecture~4 (Differentiable manifolds); \emph{Theoretische Mechanik}, Lecture~3 (Differentiable manifolds). Both introduce the $C^k$ atlas and the
differentiability of maps through charts; the classification results of this lecture (unique smooth structure up to dimension~3, uncountably many on $\R^4$) are, as far as the titles say, specific to the present course.
(From the lecture titles and the plan of the course; not checked in detail.)
\end{reading}

% ==================================================================
\section{Adding structure: refining the atlas}\label{sec:df-atlas}

Today differentiable manifolds \ts{7}{1}{0:08}. One passes from topological to differentiable manifolds by \emph{removing charts} from the topological atlas: this is the first step, ``adding structure by refining the maximal topological atlas'', the maximal $C^0$ atlas
\ts{7}{1}{0:08}\ts{7}{1}{0:34}\ts{7}{1}{0:55}.
Recall (\cref{def:mf-c0atlas,prop:mf-c0}): on a topological manifold $M_0$ there is essentially only one maximal topological atlas, because all charts one can choose are automatically $C^0$ compatible by the very definition of a topological manifold; for a topological manifold an atlas was a fully
redundant notion. Now something non-redundant is done \ts{7}{1}{0:55}\ts{7}{1}{1:16}\ts{7}{1}{1:37}.
\begin{recall}[$C^0$ compatibility]
Two charts $(U,x)$, $(V,y)$ are $C^0$ compatible if $U\cap V=\emptyset$ or the transition map $y\circ x^{-1}\colon x(U\cap V)\to y(U\cap V)$ is continuous (\cref{def:mf-compat}); for a topological manifold this always holds (\cref{prop:mf-c0}):
$x$ and $y$ are homeomorphisms, hence so is $y\circ x^{-1}$.
\end{recall}

\begin{definition}[New concepts: $C^\star$ atlas]\label{def:df-star}
Let $M_0=(M,\mathcal O)$ be a topological manifold of dimension $D$. An atlas $\mathcal A$ of $M_0$ is called a \textbf{$C^\star$ atlas} (where $C^\star$ stands for one of the compatibility classes below; the symbol is heard in the transcript as ``flower'' and written
as a placeholder, so I use $\star$ [?]) if any two charts $(U,x)$ and $(V,y)$ in $\mathcal A$ are $C^\star$ compatible, that is: either
\begin{enumerate}
  \item $U\cap V=\emptyset$ (the two chart regions, open subsets of the manifold, do not meet: then they are already compatible), or
  \item $U\cap V\ne\emptyset$ and the chart transition map
        \[
          y\circ x^{-1}\colon x(U\cap V)\to y(U\cap V)
        \]
        (both $x(U\cap V)$ and $y(U\cap V)$ are subsets of $\R^D$, $D=\dim M$) is $C^\star$ as a map $\R^D\supseteq x(U\cap V)\to y(U\cap V)\subseteq\R^D$
\end{enumerate}
\ts{7}{1}{1:37}\ts{7}{1}{2:02}\ts{7}{1}{2:28}\ts{7}{1}{2:48}\ts{7}{1}{3:09}\ts{7}{1}{3:30}\ts{7}{1}{3:51}\ts{7}{1}{4:12}.
\end{definition}
The chart transition map is not a map between manifolds but a map from $\R^{\dim M}$ to $\R^{\dim M}$. In a topological manifold it is always a homeomorphism, because $x$ and $y$ are; but now one does not want the transition maps to be pairwise $C^0$ compatible, one wants them
to be $C^\star$: ``before you think I finally turned nuts, what is $C^\star$?'' \ts{7}{1}{4:12}\ts{7}{1}{4:33}\ts{7}{1}{4:54}\ts{7}{1}{5:15}. There are several possibilities for what $C^\star$ can be:

\begin{definition}[New concepts: $C^0$, $C^k$, smooth, analytic and complex atlases]\label{def:df-ck}
For the manifold $M_0$ and an atlas $\mathcal A$:
\begin{itemize}
  \item \textbf{$C^0$:} the transition maps are continuous as maps $\R^D\to\R^D$ (as last time) \ts{7}{1}{5:15}\ts{7}{1}{5:36}.
  \item \textbf{$C^k$:} the transition maps are $k$ times continuously differentiable as maps $\R^D\to\R^D$. This is a condition pushed down to the transition maps, which is meaningful since differentiability on $\R^D$ is understood; the definition of $C^k$ \emph{for manifolds} is what
        is about to be constructed. The atlas is called a $C^k$ atlas, and one speaks of a manifold with $k$ times continuously differentiable transition functions \ts{7}{1}{5:36}\ts{7}{1}{5:57}\ts{7}{1}{6:19}\ts{7}{1}{6:40}.
  \item \textbf{$C^\infty$:} one can differentiate arbitrarily often; such atlases (and manifolds) are called \textbf{smooth}. (Very often used, as will be seen) \ts{7}{1}{6:40}\ts{7}{1}{7:01}.
  \item \textbf{$C^\omega$:} (an omega, looking a little like an infinity) \textbf{real analytic}: the transition functions are real analytic, i.e.\ can be Taylor expanded around every point of their domain; the atlas gives an \textbf{analytic manifold}
        \ts{7}{1}{7:22}\ts{7}{1}{7:43}\ts{7}{1}{8:04}.
  \item \textbf{Complex:} one speaks of a \textbf{complex atlas} and a \textbf{complex manifold} (details below) \ts{7}{1}{8:04}\ts{7}{1}{8:24}.
\end{itemize}
\end{definition}
(One argues in a second, after a few more definitions, that for $k\ge1$ there is no big difference between $C^k$ and $C^\infty$ \ts{7}{1}{7:01}\ts{7}{1}{7:22}.)
Analyticity is a stronger condition than $C^\infty$: every analytic function is $C^\infty$, but not every $C^\infty$ function is analytic, and the analysis class gives examples \ts{7}{1}{8:04}\ts{7}{1}{8:24}.

\paragraph{Complex manifolds.} So far one always mapped into some $\R^D$; how does this become complex? There are different possibilities. One can keep the definition as it is, but then one must take an even-dimensional real manifold, which is a \emph{complex} manifold if the transition functions,
being continuous anyway (``I must emphasise this, otherwise they could not come from an atlas as defined''), satisfy the \textbf{Cauchy--Riemann equations} \ts{7}{1}{8:24}\ts{7}{1}{8:46}\ts{7}{1}{9:07}\ts{7}{1}{9:27}.
\begin{definition}[Cauchy--Riemann equations]\label{def:df-cr}
Look at a map $f\colon\R^2\to\R^2$; as a set the complex plane is $\C\cong\R^2$, and a complex number is a pair $(x,y)\leftrightarrow x+iy$. Decompose the dependent variable too,
\[
  f(x+iy)=u(x,y)+i\,v(x,y),
\]
with two real functions $u,v$ of two real variables. The \textbf{Cauchy--Riemann equations} are
\[
  \frac{\partial u}{\partial x}=\frac{\partial v}{\partial y},\qquad \frac{\partial u}{\partial y}=-\frac{\partial v}{\partial x}.
\]
A \emph{continuous} map $\R^2\to\R^2$ satisfying them is as good as a complex differentiable function $\C\to\C$ \ts{7}{1}{9:27}\ts{7}{1}{9:48}\ts{7}{1}{10:08}\ts{7}{1}{10:29}\ts{7}{1}{10:49}\ts{7}{1}{11:10}\ts{7}{1}{11:32}.
\end{definition}
(A student's remark at \ts{7}{1}{12:35}: the map is differentiable from $\C$ to $\C$, not $\R^2\to\R^2$; Schuller agrees: $\R^2$ is the complex plane. The lecturer first asked whether the Cauchy--Riemann equations were known; ``no? okay, look it up''
\ts{7}{1}{9:27}\ts{7}{1}{9:48}.)
Differentiability in the complex plane is a much stronger requirement than over the reals, because very roughly one can approach the limit point in so many different ways; the field studying it is function theory (complex analysis) \ts{7}{1}{11:32}\ts{7}{1}{11:52}. With more variables one
gets indices, $\partial u^a/\partial x^b=\ldots$, and one imposes the Cauchy--Riemann equations on the transition maps; since one always has to split into two, this requires $\dim M$ to be even as a real manifold \ts{7}{1}{11:52}\ts{7}{1}{12:13}\ts{7}{1}{12:35}.
\begin{coursenote}[The higher-dimensional equations]
For $z^a=x^a+iy^a$, $a=1,\dots,n$ and a map $f^a=u^a+iv^a$ the equations read $\partial u^a/\partial x^b=\partial v^a/\partial y^b$, $\partial u^a/\partial y^b=-\partial v^a/\partial x^b$; the dimension of the real manifold is then $2n$. Not written out in class (``you get indices here'').
\end{coursenote}
So one can add additional or different structure and pick which one one wants; by far the most important are the smooth differentiable structures and the complex ones, for reasons explained in a moment \ts{7}{1}{12:35}\ts{7}{1}{12:55}.

% ==================================================================
\section{Smooth structure is not a big choice}\label{sec:df-whitney}

\begin{theorem}[Whitney]\label{thm:df-whitney}
Let $k\ge1$ (very important: this does \emph{not} include $k=0$). Any maximal $C^k$ atlas contains a $C^\infty$ atlas; and two maximal $C^k$ atlases that contain the same $C^\infty$ atlas are already identical
\ts{7}{1}{12:55}\ts{7}{1}{13:17}\ts{7}{1}{13:39}\ts{7}{1}{14:01}\ts{7}{1}{14:26}.
\end{theorem}
(The theorem is stated, not proved, in class; the name is given as ``I think Whitney''.)
\begin{coursenote}[On the attribution and the wording]
In the transcript the name is ``Whitney'' with ``I think'', and the statement is a loose one; the usual form, due to H.~Whitney (1936), is that every $C^k$ structure, $k\ge1$, contains a compatible smooth one, unique up to diffeomorphism.
Not said in class; ``contains'' is as in class.
\end{coursenote}
\emph{Implication.} Once one has found a $C^1$ atlas for a manifold, it is not a stronger condition to say that the manifold is already $C^\infty$: it is already smooth; a $C^1$ atlas can be turned into a $C^\infty$ one by removing more charts \ts{7}{1}{14:26}\ts{7}{1}{14:49}\ts{7}{1}{15:10}.
The converse is the content of the next sentence: one could have a topological manifold with \emph{no} differentiable atlas at all; but once one sees one, one can go all the way down to $C^\infty$ \ts{7}{1}{15:10}\ts{7}{1}{15:30}. So it may be clever to say that a theorem valid for $C^2$ manifolds is automatically valid for
$C^3,C^4,C^5,\dots$ since they are contained, and one can construct from a $C^2$ (even a $C^1$) manifold a smooth one. When talking about differentiable manifolds or the existence of a differentiable structure one need not distinguish between the cases $C^k$ with $k\ge1$ and $C^\infty$, so
smooth in the above sense \ts{7}{1}{15:30}\ts{7}{1}{15:51}\ts{7}{1}{16:11}\ts{7}{1}{16:37}.

% ==================================================================
\section{Differentiable manifolds}\label{sec:df-manifold}

What is done with the differentiable structure? First one defines the notion of a differentiable manifold, then differentiable maps between two manifolds; these are the structure preserving maps if they are also invertible, which gives a \emph{classification} of smooth manifolds, with the first surprises, even
from the point of view of fundamental physics \ts{7}{1}{16:37}\ts{7}{1}{17:00}\ts{7}{1}{17:22}.

\begin{definition}[New concepts: $C^k$ manifold]\label{def:df-manifold}
A \textbf{$C^k$ manifold} is a triple $(M,\mathcal O,\mathcal A)$ where $(M,\mathcal O)$ is a topological manifold (``of course that stays'') and $\mathcal A$ is a \emph{maximal} $C^k$ atlas \ts{7}{1}{17:22}\ts{7}{1}{17:45}\ts{7}{1}{18:09}.
\end{definition}
One has a new additional layer of structure: the $C^k$ atlas is always a subset of a $C^0$ atlas, one takes out strategically certain charts \ts{7}{1}{18:09}\ts{7}{1}{18:31}.

\begin{remark}[Incompatible atlases]\label{rem:df-incompat}
A given topological manifold can carry different \emph{incompatible} atlases. Only compatibility of charts was defined, but what is meant is clear: technically two atlases are compatible if their union is again an atlas of the same type; or: if between any two charts of the two atlases the
compatibility of type $C^\star$ holds \ts{7}{1}{18:31}\ts{7}{1}{18:58}\ts{7}{1}{19:18}.
\end{remark}

\begin{example}[Two incompatible smooth atlases of the real line]\label{ex:df-line}
Let the underlying topological manifold be $\R$ with the standard topology \ts{7}{1}{19:18}\ts{7}{1}{19:39}.
\begin{itemize}
  \item \emph{Atlas one:} $\mathcal A_1=\{(\R,\mathrm{id}_\R)\}$, a single chart covering all of $\R$ (as before, \cref{ex:mf-atlases}). It is trivially a $C^\infty$ atlas, even $C^1$, everything: the chart overlaps only with itself and $\mathrm{id}\circ\mathrm{id}^{-1}=\mathrm{id}$ is $C^\infty$
        \ts{7}{1}{19:39}\ts{7}{1}{20:00}\ts{7}{1}{20:21}.
  \item \emph{Atlas two:} $\mathcal A_2=\{(\R,x)\}$ again with the entire real line as chart region, but with chart map
        \[
          x\colon\R\to\R,\qquad a\mapsto a^{1/3}
        \]
        (the third root) \ts{7}{1}{20:21}\ts{7}{1}{20:41}. First check that this is a chart map at all: it is one-to-one, continuous, invertible with the third power as inverse, which is continuous too: a homeomorphism from the manifold $\R$ down to the chart $\R^1$. Its transition map
        with itself, $x\circ x^{-1}=\mathrm{id}$, is $C^\infty$, so this is also a $C^\infty$ atlas \ts{7}{1}{20:41}\ts{7}{1}{21:03}\ts{7}{1}{21:25}\ts{7}{1}{21:46}\ts{7}{1}{22:06}.
\end{itemize}
Both can be made into \emph{maximal} $C^\infty$ atlases by adding all the other charts whose transition functions with the chart(s) already there are $C^\infty$ (call them $\mathcal A_1^{\max}$ and $\mathcal A_2^{\max}$) \ts{7}{1}{22:06}\ts{7}{1}{22:27}.
But the two atlases are \emph{incompatible}: were they compatible, their union would be another $C^\infty$ atlas, but the transition map from $(\R,\mathrm{id})$ to $(\R,x)$, $x\circ\mathrm{id}^{-1}=a\mapsto a^{1/3}$, though continuous (it must be), is not $C^\infty$, not even $C^1$: its derivative at $0$ is infinite, not defined
\ts{7}{1}{22:27}\ts{7}{1}{22:48}\ts{7}{1}{23:10}\ts{7}{1}{23:31}\ts{7}{1}{23:53}. So the two charts are not $C^1$ compatible: as long as they are parts of different atlases that is no problem, but they could not be put into the same atlas and give a $C^\infty$ atlas \ts{7}{1}{23:31}\ts{7}{1}{23:53}.
One can therefore equip the real line with at least two different incompatible $C^\infty$ structures (just one example) \ts{7}{1}{23:53}\ts{7}{1}{24:13}.
\end{example}
That looks bad: one seems to have to decide which smooth structure to establish; it is not so much a choice (one would have to do it for all the $C^k$ levels, which would be even worse), but even looking only at smooth structures, a topological manifold does not have a unique one. If one wants to do physics, and physics
wants differentiability, ``what a disaster this is: not even the real line has a unique differentiable structure'' \ts{7}{1}{24:13}\ts{7}{1}{24:34}\ts{7}{1}{24:56}. The situation is not so bad, but to explain this one needs the next definition \ts{7}{1}{24:56}\ts{7}{1}{25:17}.

% ==================================================================
\section{Differentiable maps and diffeomorphisms}\label{sec:df-maps}

A given $C^k$ manifold is equipped with one specific choice of atlas; assume this pick is made (which pick? \ts{7}{1}{25:17}) \ts{7}{1}{25:17}\ts{7}{1}{25:38}.
\begin{definition}[New concepts: differentiable map between manifolds]\label{def:df-map}
Let $\phi\colon M\to N$ be a map, where $(M,\mathcal O_M,\mathcal A_M)$ and $(N,\mathcal O_N,\mathcal A_N)$ are $C^k$ manifolds (for a map it would suffice that $M,N$ are sets; for a continuous map topological spaces would be needed; now both are differentiable manifolds)
\ts{7}{1}{25:17}\ts{7}{1}{25:38}\ts{7}{1}{25:59}. Then $\phi$ is called \textbf{differentiable at the point $p\in M$} (of class $C^k$ at $p$) if for \emph{some} chart $(U,x)$ in the atlas $\mathcal A_M$ with $p\in U$ and \emph{some} chart $(V,y)$ in the atlas $\mathcal A_N$ with $\phi(p)\in V$ the map
\[
  y\circ\phi\circ x^{-1}\colon\ \R^{\dim M}\supseteq x(U)\longrightarrow y(V)\subseteq\R^{\dim N}
\]
is $C^k$ as a map from $\R^{\dim M}$ to $\R^{\dim N}$ \ts{7}{1}{25:59}\ts{7}{1}{26:24}\ts{7}{1}{26:45}\ts{7}{1}{27:10}\ts{7}{1}{27:32}\ts{7}{1}{27:52}\ts{7}{1}{28:13}\ts{7}{1}{28:39}\ts{7}{1}{29:00}.
\end{definition}
(One can only define differentiability of a map if a differentiable atlas has been chosen: for topological manifolds or spaces there is no such notion. The definition uses the chosen atlas in an essential way; which one does not matter, one chooses one \ts{7}{1}{26:24}\ts{7}{1}{26:45}.)
Why is $y\circ\phi\circ x^{-1}$ a map $\R^{\dim M}\to\R^{\dim N}$? Look at $M$, but only at the region $U$ that contains the point of interest, and at $N$, but only at the region $V$ containing the image of that point under $\phi$; the chart maps bring one down to $x(U)\subseteq\R^{\dim M}$
and $y(V)\subseteq\R^{\dim N}$ ($\dim M$ and $\dim N$ may differ, they are different manifolds). One wants to speak of a map between differentiable manifolds, but because an atlas is available one does so \emph{only} through the available atlases: pick charts that cover the point of interest and its
image, and consider the map in its \emph{chart representation}: mimic $\phi$ by the path $x^{-1}$, then $\phi$, then $y$. This is a map from $\R^{\dim M}$ to $\R^{\dim N}$, and for these one knows what differentiability means; so the differentiability of $\phi$ is \emph{defined} by that of its chart representation
\ts{7}{1}{29:00}\ts{7}{1}{29:21}\ts{7}{1}{29:42}\ts{7}{1}{30:02}\ts{7}{1}{30:22}\ts{7}{1}{30:43}\ts{7}{1}{31:04}.

The definition says ``for some chart''; one must check that the notion does not depend on the charts chosen: well-definedness \ts{7}{1}{31:04}\ts{7}{1}{31:25}\ts{7}{1}{31:46}.
\begin{proposition}[Well-definedness]\label{prop:df-welldef}
The notion of differentiability at $p$ of \cref{def:df-map} does not depend on the chosen charts: if $y\circ\phi\circ x^{-1}$ is $C^k$ for one choice of charts $(U,x)\in\mathcal A_M$, $(V,y)\in\mathcal A_N$, then $\tilde y\circ\phi\circ\tilde x^{-1}$ is $C^k$ for every other choice
$(\tilde U,\tilde x)\in\mathcal A_M$, $(\tilde V,\tilde y)\in\mathcal A_N$ with $p\in\tilde U$ and $\phi(p)\in\tilde V$ \ts{7}{1}{31:46}\ts{7}{1}{32:11}.
\end{proposition}
\begin{proof}
Take $U,V$ the subsets of $M,N$ containing $p$ and $\phi(p)$, with charts $x\colon U\to x(U)$ and $y\colon V\to y(V)$ from the respective atlases (very important: from the $C^k$ atlases), and the map $y\circ\phi\circ x^{-1}$ that is $C^k$ as a map $\R^{\dim M}\to\R^{\dim N}$
\ts{7}{1}{32:33}\ts{7}{1}{32:53}\ts{7}{1}{33:15}\ts{7}{1}{33:35}\ts{7}{1}{33:56}. Now let different charts $\tilde x$, $\tilde y$ be chosen, being extra picky: they have their own chart domains $\tilde U\ni p$, $\tilde V\ni\phi(p)$; replacing $U$ by $U\cap\tilde U$ and $V$ by $V\cap\tilde V$ (``otherwise the picture gets
messy'', without loss of generality) one may assume the domains are the same. One considers the map $\tilde y\circ\phi\circ\tilde x^{-1}$; a priori nothing guarantees that it is also $C^k$ \ts{7}{1}{33:56}\ts{7}{1}{34:16}\ts{7}{1}{34:36}\ts{7}{1}{34:57}\ts{7}{1}{35:17}\ts{7}{1}{35:38}.
But it can be written as a composition
\[
  \tilde y\circ\phi\circ\tilde x^{-1}=\bigl(\tilde y\circ y^{-1}\bigr)\circ\bigl(y\circ\phi\circ x^{-1}\bigr)\circ\bigl(x\circ\tilde x^{-1}\bigr),
\]
reading off the diagram \ts{7}{1}{35:38}\ts{7}{1}{35:58}. Here $\tilde y\circ y^{-1}$ is a chart transition map between two charts of the same $C^k$ atlas $\mathcal A_N$, hence $C^k$ as a map $\R^{\dim N}\to\R^{\dim N}$ (twice the same dimension, because it is a transition map: for $M$ and $N$ the dimensions could differ), and
$x\circ\tilde x^{-1}$ is a chart transition map between two charts of $\mathcal A_M$, hence $C^k$ as a map $\R^{\dim M}\to\R^{\dim M}$ \ts{7}{1}{35:58}\ts{7}{1}{36:18}\ts{7}{1}{36:39}\ts{7}{1}{37:00}\ts{7}{1}{37:20}.
The middle factor is $C^k$ by assumption, and a composition of $C^k$ maps (over $\R^n$, $\R^m$ and so on) is again $C^k$; hence once it is checked in one chart it is true in any other chart \ts{7}{1}{37:20}\ts{7}{1}{37:41}\ts{7}{1}{38:02}.
\end{proof}
Now one sees why it was so important to restrict the charts to those of a $C^k$ atlas: otherwise the compatibility condition would not be there; if the transition maps were only $C^0$, as they are for any topological manifold, one could check $C^k$ in one chart, but $C^k$ composed with $C^0$ maps
possibly destroys the $C^k$ property and leaves a merely continuous map. The chart transition maps must be at least of the same differentiability class: hence \emph{on a $C^k$ manifold one can define $C^k$ differentiability for maps between $C^k$ manifolds, but no higher}, because it might be
destroyed \ts{7}{1}{38:02}\ts{7}{1}{38:23}\ts{7}{1}{38:44}\ts{7}{1}{39:06}.

\begin{intuition}[The manifold philosophy]\label{int:df-philosophy}
You want to define something for an object on the manifold, and you define it by looking at its representative in a chart, or in several charts (here one goes from one manifold to the other, so two charts at least); then you show that it is independent of the choice of chart, and for
well-definedness to work you must suitably restrict the transition functions. For a complex manifold, one could check complex differentiability in one representation, but if the transition maps did not satisfy the higher-dimensional Cauchy--Riemann equations on top of being continuous anyway, one could not conclude that what was
checked in one chart is true in any other, hence could not attribute it as a property to the abstract object. For topological manifolds the chart picture was redundant, since continuity can be defined in its own right between topological spaces; once one goes beyond continuity, the chart picture becomes essential: one
steals from the knowledge of differentiability on $\R^n\to\R^m$ and lifts the notion to the manifold level \ts{7}{1}{39:06}\ts{7}{1}{39:26}\ts{7}{1}{39:47}\ts{7}{1}{40:08}\ts{7}{1}{40:30}\ts{7}{1}{40:54}.
\end{intuition}

\begin{definition}[New concepts: diffeomorphism, diffeomorphic manifolds]\label{def:df-diffeo}
A map $\phi\colon M\to N$ between smooth manifolds is called a \textbf{diffeomorphism} if it is bijective as a map between sets and both $\phi$ and its inverse $\phi^{-1}\colon N\to M$ are $C^\infty$ (one could say $C^k$, but let us say $C^\infty$)
\ts{7}{1}{40:54}\ts{7}{1}{41:16}\ts{7}{1}{41:38}. Two smooth manifolds $(M,\mathcal O_M,\mathcal A_M)$ and $(N,\mathcal O_N,\mathcal A_N)$ are called \textbf{diffeomorphic}, written $M\simeq_{\mathrm{diff}}N$, if there exists a diffeomorphism between them
\ts{7}{1}{41:38}\ts{7}{1}{42:00}\ts{7}{1}{42:22}.
\end{definition}
The diffeomorphisms are the isomorphisms, the structure preserving maps between differentiable (smooth) manifolds: ``always the same pattern'' (\cref{int:cl-theme}). Strictly, one should write that the triple $(M,\mathcal O_M,\mathcal A_M)$ is diffeomorphic to the triple $(N,\mathcal O_N,\mathcal A_N)$, as that is the full name of a smooth manifold:
a set, a topological structure and a pick of a smooth atlas. In future the notation is lightened and only the sets $M$ and $N$ are written, but ``if I say these $M$ and $N$ be smooth manifolds you know that secretly they carry with them a choice of topology and a choice of atlas'', and another manifold carries its own possibly
different choices; this must not be suppressed in the mind \ts{7}{1}{42:22}\ts{7}{1}{42:42}\ts{7}{1}{43:03}\ts{7}{1}{43:23}. Being diffeomorphic is an equivalence relation (every manifold is diffeomorphic to itself, as the identity is a diffeomorphism) \ts{7}{1}{43:23}\ts{7}{1}{43:44}, and it is custom not to distinguish
diffeomorphic manifolds: they are called the same, seen through the glasses of differentiable structures \ts{7}{1}{43:44}\ts{7}{1}{44:04}\ts{7}{1}{44:25}.

\begin{exercise}[The two atlases of the real line are diffeomorphic]\label{ex:df-line-diffeo}
\emph{Not said in class.} Show that $(\R,\mathcal A_1^{\max})$ and $(\R,\mathcal A_2^{\max})$ of \cref{ex:df-line} are diffeomorphic: take $\phi\colon(\R,\mathcal A_2^{\max})\to(\R,\mathcal A_1^{\max})$, $\phi(p)=p^{1/3}$, and compute its chart representation with the charts $x=(\cdot)^{1/3}$ and $\mathrm{id}$,
and likewise for $\phi^{-1}$. This anticipates \cref{sec:df-classification}.
\end{exercise}

% ==================================================================
\section{Levels of structure}\label{sec:df-levels}

Remark: it is custom to consider diffeomorphic smooth manifolds to be the same, and this is all a question of the \emph{level of structure} \ts{7}{1}{44:04}\ts{7}{1}{44:25}.
Two sets $M$ and $N$ may be the same seen through the glasses of set theory (there is a bijection); equip each with a topology, and unless there exists a homeomorphism between the extended structures they are no longer the same seen through the glasses of topology. One can push this on: even if they are
homeomorphic as topological spaces, one can equip them with different differentiable structures, and then they may be different as differentiable manifolds. This was seen with the real line with the standard topology: trivially homeomorphic to itself, but with one smooth atlas on one copy and another on the other
copy they are not the same as differentiable manifolds (\cref{ex:df-line}; but see below, they are in fact diffeomorphic) \ts{7}{1}{44:25}\ts{7}{1}{44:46}\ts{7}{1}{45:07}\ts{7}{1}{45:27}\ts{7}{1}{45:47}.
\begin{coursenote}[A slip in the example]
The lecture says that with two different atlases the two copies of the real line are ``no longer the same as differentiable manifolds''; what is true is that the atlases are \emph{incompatible}, i.e.\ different structures on the same set; whether the resulting manifolds are diffeomorphic is
a different question, answered in the next section (they are diffeomorphic, \cref{ex:df-line-diffeo}). Not said in class.
\end{coursenote}
``As soon as I start painting people's hair red and green they are no longer the same''; it always depends on context what ``these manifolds are the same'' means \ts{7}{1}{45:47}\ts{7}{1}{46:09}.

\paragraph{A student asks about group structures.} A student: one had said that some sets are equal or isomorphic as a group structure; is it a linear stacking? Schuller goes back to page one (the picture of the first lecture): a set can be equipped with a topology (modify the picture), or instead of a topology
with a map $\diamond\colon M\times M\to M$ (``some blob or diamond''); if $\diamond$ is associative, has a neutral element in $M$ and inverse elements, it is a \textbf{group}. A set can be equipped with a topology \emph{and} such a structure: one obtains a \textbf{topological group}, also called a Lie group
(see the coursenote). On top of the topology one can say it is a manifold and then put a differentiable structure: a differentiable manifold, which is the path followed so far \ts{7}{1}{46:09}\ts{7}{1}{46:29}\ts{7}{1}{46:49}\ts{7}{1}{47:12}\ts{7}{1}{47:32}\ts{7}{1}{47:54}.
\begin{coursenote}[Topological versus Lie groups]
The lecture says that a set with a topology and a group structure is ``a topological group, also called a Lie group'', which is a slip: a Lie group is a group that is also a \emph{smooth manifold}, with smooth group operations (\cref{ch:l13}). What is true, and is the content of the following paragraph, is that a
locally Euclidean topological group is already a Lie group. Not said in class.
\end{coursenote}
\emph{Hilbert's problem.} Once there, with a set with a topology and a group structure, one has a so-called Lie group; one can continue. Around 1900 Hilbert gave a big speech and said that there are 22 problems to be solved in the twentieth century; one of them was whether every topological group (called Lie group)
is already a differentiable Lie group, and \emph{the answer is yes}. Hence the analysis of Lie groups is best done by already employing a differentiable structure, as in the course: the differentiable structure gives rise to the so-called Lie algebras, and by Lie algebras one can study Lie groups very efficiently, at least some
aspects \ts{7}{1}{47:54}\ts{7}{1}{48:14}\ts{7}{1}{48:35}\ts{7}{1}{48:55}.
\begin{history}[Hilbert's fifth problem]
\emph{Checked.} At the Paris congress of 1900 Hilbert listed 23 problems (published version; the lecture says 22, as heard); the fifth asks whether a continuous group is automatically a differentiable one. It was solved in the affirmative, for locally Euclidean groups, by Gleason and by
Montgomery and Zippin in 1952.\par
\emph{References.} D.~Hilbert, ``Mathematische Probleme'', \emph{G\"ottinger Nachrichten} (1900); A.~Gleason, \emph{Ann.\ Math.}\ 56 (1952); D.~Montgomery, L.~Zippin, \emph{Ann.\ Math.}\ 56 (1952).
\end{history}
So what this picture should convey is that one can have one structure and not the other, but not a differentiable structure without a topological space, and not a topological space without a set; but one can combine the structures in various ways. Another thing: an additional structure like a scalar multiplication
$S\colon K\times M\to M$ (``an S multiplication''), and if the right eight axioms hold, it is a \textbf{vector space}. A vector space can be equipped with a topology and so on; ``you can combine the structures in almost arbitrary ways, but not completely arbitrary ways, because some structures require the other structure to underlie''.
It is always very important whether one talks about a vector space or about a topological vector space \ts{7}{1}{48:55}\ts{7}{1}{49:15}\ts{7}{1}{49:36}\ts{7}{1}{49:57}\ts{7}{1}{50:18}.
\begin{secondary}[Not a linear stacking of structures]
A second student question: for the fundamental groups one says the groups are isomorphic as groups, while the torus and the cylinder are not isomorphic topologically; could one then say that they cannot be isomorphic in the group sense? No: the picture talks about \emph{the same set} being equipped with more
and more structure of this and that type, in that combination. If one looks at a topological space, one can construct from it the fundamental group (\cref{sec:iv-pi1}); but the group has its own underlying set (for the torus, $\mathbb Z\times\mathbb Z$, a different set) which carries a group structure, and one can build another
``building like this'' on that set. It does not mean that in one mathematics problem things always come in this one linear stacking; the picture, which ``looks pretty, maybe more confusing than enlightening'', was not meant to convey that \ts{7}{1}{50:18}\ts{7}{1}{50:38}\ts{7}{1}{50:59}\ts{7}{1}{51:19}\ts{7}{1}{51:40}\ts{7}{1}{52:00}.
\end{secondary}

% ==================================================================
\section{How many smooth structures? The classification}\label{sec:df-classification}

Now that diffeomorphic manifolds are essentially the same, return to the question \ts{7}{1}{52:00}\ts{7}{1}{52:21}: how many differentiable structures are there up to isomorphism? For instance, is the real line with the standard topology and the maximal extension of the atlas $\mathcal A_1$ of \cref{ex:df-line} at
least diffeomorphic to the real line with the standard topology and (the maximal extension of) $\mathcal A_2$? They were clearly not the same atlas, $\mathcal A_1^{\max}\ne\mathcal A_2^{\max}$: one chart in each, but once one has a chart one adds all compatible charts and gets the maximisation; and one started worrying which
structure one would use as \emph{the} differentiable structure \ts{7}{1}{52:21}\ts{7}{1}{52:43}\ts{7}{1}{53:06}\ts{7}{1}{53:27}. But one should not be so precise and look only at the atlas; one should ask whether the resulting differentiable manifold differs \emph{essentially}: a non-essential difference
is one for which they are diffeomorphic. Then from the point of view of differentiability they are the same, since every function differentiable on one can be pushed over to the other and is differentiable there \ts{7}{1}{53:27}\ts{7}{1}{53:47}\ts{7}{1}{54:08}.
There are rather deep results, not proved here, some of them big achievements. The question:
\begin{center}\emph{How many different differentiable structures can one put on a given topological manifold, up to diffeomorphism?}\end{center}
(I hand you a topological manifold, you have to decide which differentiable structure to put on it) \ts{7}{1}{54:08}\ts{7}{1}{54:28}\ts{7}{1}{54:49}\ts{7}{1}{55:10}. The surprising answer depends on the dimension \ts{7}{1}{55:10}\ts{7}{1}{55:31}.

\begin{theorem}[Dimensions $1,2,3$ (Rad\'o--Moise)]\label{thm:df-low}
For $\dim M=1,2$ or $3$ every topological manifold carries, up to diffeomorphism, a \emph{unique} smooth manifold structure \ts{7}{1}{55:31}\ts{7}{1}{55:52}\ts{7}{1}{56:12}\ts{7}{1}{56:34}. (In the transcript: ``Radon Mois theorems''.)
\end{theorem}
Although one could choose different smooth charts already for the real line, and the atlases are not compatible, the resulting manifolds are diffeomorphic; so everything one does with differentiability is unique. One now has a notion of differentiability on the torus, say: ``see what an achievement this is''. At high school one learns
differential calculus of one variable over the real line, at university in an engineering course over $\R^n$ or $\R^3$, and now one can do it on arbitrary manifolds, with holes and whatnot; up to dimension three everything is unique \ts{7}{1}{56:34}\ts{7}{1}{56:58}\ts{7}{1}{57:19}\ts{7}{1}{57:39}.

\begin{theorem}[Dimension greater than four: finitely many]\label{thm:df-high}
For $\dim M>4$ there are only finitely many different smooth manifolds (up to diffeomorphism) that can be made of a given topological manifold \ts{7}{1}{58:01}\ts{7}{1}{1:00:04}\ts{7}{1}{1:00:24}\ts{7}{1}{1:00:46}.
\end{theorem}
Dimension four is left out for now: ``dimension of $M$ greater than four, but \emph{not} four: decidedly not four'' \ts{7}{1}{57:39}\ts{7}{1}{58:01}. There is a toolbox called \emph{surgery theory}: take, say, a sphere and operate on it like a surgeon: cut out two pieces with a scalpel; a colleague gives
a cylinder, which is pushed in and sewn at both ends; what one gets is the torus. The idea (a big field): the sphere and the cylinder are well understood, and if one performs surgery in such a way that one controls the topological invariants one cares about (the fundamental group, later homology), one understands the torus;
that is a trivial example, the aim being to understand all manifolds of dimension higher than four by reducing them to elements and surgery procedures one understands. Done systematically, in the 1960s, this gave the theorem \ts{7}{1}{58:01}\ts{7}{1}{58:22}\ts{7}{1}{58:43}\ts{7}{1}{59:03}\ts{7}{1}{59:23}\ts{7}{1}{59:44}\ts{7}{1}{1:00:04}.
(In the lecture it is first said ``dim $M\ge4$'' and then ``greater than four''; see the coursenote.)
\begin{coursenote}[Dimension four and the wording]
The lecture first says ``greater or equal to four'' and then ``greater than four'' and ``not four''; the correct statement is $\dim M\ge5$, with $\dim M=4$ treated separately, as done next. Not said in class.
\end{coursenote}
That is not as neat as dimensions one, two, three, but the structures are finitely many, so there is a chance, in principle, to enumerate them and write them all down. Take say the four-sphere (``let us find all the differentiable structures''; for a given manifold it may be very difficult in practice, but the theorem
says that they are finitely many) \ts{7}{1}{1:00:46}\ts{7}{1}{1:01:07}.
\begin{application}[If space-time were higher dimensional]\label{app:df-physics}
Some physicists believe that space-time is, say, 10-dimensional (``not proposing or condemning that''); then one must worry which smooth structure to equip the manifold with. Physics can make experiments: if different structures lead to different observations, one can at least enumerate them and make finitely many
observations to distinguish among them, in principle. So higher dimension is in a sense good news: one could determine by experiment which of the structures nature has chosen, if any (if space-time is a differentiable manifold in the first place). ``This has immediate impact on physics, if you take it seriously, and we want to do that''
\ts{7}{1}{1:01:07}\ts{7}{1}{1:01:29}\ts{7}{1}{1:01:49}\ts{7}{1}{1:02:10}\ts{7}{1}{1:02:30}.
\end{application}

\begin{theorem}[Dimension four]\label{thm:df-four}
A $4$-dimensional smooth manifold, like $\R^4$: up to diffeomorphism there are \emph{uncountably many} different smooth manifolds that can be made of the same $4$-dimensional topological manifold; in particular for $\R^4$ with the standard topology, a topological manifold of dimension four.
This holds at least in the non-compact case (c.1)
\ts{7}{1}{1:02:30}\ts{7}{1}{1:02:51}\ts{7}{1}{1:03:14}\ts{7}{1}{1:03:37}\ts{7}{1}{1:03:58}.
\end{theorem}
The \emph{compact case} (c.2) has only partial results: if the second \textbf{Betti number} $b_2$, a topological invariant, is greater than $18$ [?: transcript ``greater than 18''], then there are only (countably) infinitely many different smooth manifolds; not much better, so the compact case is a little
better than the non-compact one but still infinite \ts{7}{1}{1:03:58}\ts{7}{1}{1:04:21}\ts{7}{1}{1:04:45}\ts{7}{1}{1:05:05}\ts{7}{1}{1:05:25}\ts{7}{1}{1:05:46}.
\begin{coursenote}[Which compact case]
The transcript has ``countably infinitely many'' for $b_2>18$; the lecture's statement is only a partial result and I could not check the exact hypothesis; marked [?]. Not said in class.
\end{coursenote}
\begin{definition}[New concepts: Betti numbers (as explained informally)]\label{def:df-betti}
The Betti numbers are defined in terms of homology groups (to come later); intuitively the zeroth Betti number $b_0$ is the number of connected pieces of the space, $b_1$ the number of circular (one-dimensional) holes, $b_2$ the number of two-dimensional holes, and so on; ``the pretzel but with 19 holes or so'' has a large number of holes
\ts{7}{1}{1:04:21}\ts{7}{1}{1:04:45}\ts{7}{1}{1:05:05}\ts{7}{1}{1:05:25}. (Name: after E.~Betti, who introduced the numbers in 1871; checked.)
\end{definition}
\begin{application}[Special relativity on $\R^4$]\label{app:df-r4}
When one says space-time is $\R^4$ with some additional structure (the light cone or the metric, in special relativity) and one finds tangent vectors to world lines, the question is which smooth structure to put on $\R^4$ as a topological manifold in order to do special relativity: there are uncountably many to choose
from. Can one do experiments to distinguish between them? One would need, with luck, uncountably many experiments to find the one: ``it's not so good, right? But it's a result''. ``Unfortunately I have nothing better to say about this, it's just an observation, and this is what the analysis yields''. From a philosophical point of view it is
serious (``I do not know whether from a scientific practical point of view''): in principle our theories could fail on $\R^4$ because the wrong differentiable structure was chosen, and might be repaired by choosing another; ``I don't think this is a point where you could start doing research, but maybe you can. As a mathematical result it stays
and should be noted'' \ts{7}{1}{1:05:46}\ts{7}{1}{1:06:07}\ts{7}{1}{1:06:29}\ts{7}{1}{1:06:53}\ts{7}{1}{1:07:14}\ts{7}{1}{1:07:34}.
\end{application}
\begin{history}[Exotic spheres and exotic $\R^4$]
\emph{Not said in class.} \emph{Checked.} J.~Milnor found in 1956 smooth structures on the 7-sphere that are not diffeomorphic to the standard one (exotic spheres); for $\R^4$, exotic structures follow from the work of Donaldson and Freedman (1982), and C.~Taubes (1987) showed that there are uncountably many
non-diffeomorphic ones.\par
\emph{References.} J.~Milnor, \emph{Ann.\ Math.}\ 64 (1956); S.~Donaldson, \emph{J.\ Diff.\ Geom.}\ 18 (1983); M.~Freedman, \emph{J.\ Diff.\ Geom.}\ 17 (1982); C.~Taubes, \emph{J.\ Diff.\ Geom.}\ 25 (1987).
\end{history}

% ==================================================================
\section{Outlook: tangent spaces without embedding}\label{sec:df-outlook}

So a differentiable structure is defined, one can classify; but what is really new about differentiable manifolds is that, roughly speaking, \emph{they have tangent spaces} \ts{7}{1}{1:07:14}\ts{7}{1}{1:07:34}\ts{7}{1}{1:07:54}. The key feature: at each point of the manifold there exists a tangent space (the words are in quotes: one has an idea of what a tangent space looks like,
but that idea will not be used; it will be done in a different way) \ts{7}{1}{1:07:54}\ts{7}{1}{1:08:16}\ts{7}{1}{1:08:36}.

\begin{intuition}[Why not a plane in the ambient space]\label{int:df-embedding}
Look at the sphere $S^2$ as a topological manifold; one always thinks of it as embedded in three-dimensional space, but as a topological manifold it does not care: the set and the topology make no reference to what lies outside the surface; the same holds when one adds a smooth atlas. Now explain tangent spaces in the
following way: take one point of the sphere, rotate it ``for dramatic purposes'' to the horizon, and imagine the tangent space as all points or vectors in the tangent plane. This makes reference to the space in which the sphere is embedded, which has nothing to do with the smooth structure, independent of the embedding. So tangent spaces
will \emph{not} be defined as planes in an exterior space in which the manifold is embedded. (Due to various theorems, one could do it that way without loss of generality \ts{7}{1}{1:08:36}\ts{7}{1}{1:08:56}\ts{7}{1}{1:09:17}\ts{7}{1}{1:09:37}\ts{7}{1}{1:09:58}\ts{7}{1}{1:10:19}\ts{7}{1}{1:10:39}.)
Nevertheless it is not done, because if the manifold is not two- but four-dimensional and represents our universe, space-time, it would be funny to think of the universe as embedded, physically speaking, in a bigger space in which the tangent spaces lie, since that bigger space would obviously be what should be called the universe,
the thing that comprehends everything. Even from the point of view of theory building in physics it is nice to construct everything \emph{intrinsically}, without recourse to structures outside the object: one must define tangent spaces solely with the structure one is handed, notwithstanding intuition (``you may think of tangent spaces
this way, but it is better not to define them this way'') \ts{7}{1}{1:10:39}\ts{7}{1}{1:11:01}\ts{7}{1}{1:11:22}\ts{7}{1}{1:11:42}\ts{7}{1}{1:12:03}.
\end{intuition}

In the next subsections one fixes a point of the manifold and constructs the tangent space at this point, its dual, and the tensor spaces over this point; later all tangent spaces are looked at together: tangent vectors at one point form a vector space, but taking at every point a different vector gives a
\emph{vector field} on the manifold \ts{7}{1}{1:12:03}\ts{7}{1}{1:12:23}\ts{7}{1}{1:12:44}. Vector spaces are linear algebra: a few linear algebra notions will be reviewed quickly, in particular the construction of tensors, because ``there are some standard confusions that we will now try to eliminate once and for all''
\ts{7}{1}{1:12:44}\ts{7}{1}{1:13:04}.
However, vector fields still have a linear structure: two vector fields can be added, but also scaled, and at different points of the manifold one can scale differently, with a \emph{function}. It turns out that the strict notion of a vector space is too restrictive to deal in a clever way with vector fields, and one goes to another
structure called a \textbf{module}. A show of hands (``who of you saw modules in linear algebra? Oh, fantastic'' --- some had): it will be new for most, ``it's my responsibility, all your future misconceptions will be my responsibility''; normally one meets modules in linear algebra two, a very useful structure, which also in physics
``we use all the time, we just choose to ignore that they are not really vector spaces anymore'' \ts{7}{1}{1:13:04}\ts{7}{1}{1:13:25}\ts{7}{1}{1:13:46}\ts{7}{1}{1:14:07}. This is the plan of the next two lectures: tensor space theory over a field (\cref{ch:l08}) and then over a ring.

% ==================================================================
\section*{Key concepts}
\begin{keyconcepts}
\textbf{Retrieval questions.}
\begin{enumerate}
  \item What does it mean for an atlas to be $C^k$, $C^\infty$, $C^\omega$, complex? Which condition do the transition maps satisfy in each case?
  \item State the Whitney theorem on maximal $C^k$ atlases. Why does it allow one to speak of ``smooth'' only? Why must $k\ge1$?
  \item Define a $C^k$ manifold. Why can a topological manifold carry incompatible atlases? Give the example of $\R$ with the atlases of the identity and of $a\mapsto a^{1/3}$.
  \item Define the differentiability of a map between manifolds at a point. Prove that it does not depend on the charts chosen: where does the compatibility of the atlases enter?
  \item Why can one define $C^k$ differentiability on a $C^k$ manifold but no higher?
  \item Define diffeomorphism and diffeomorphic manifolds. Why is the real line with the two atlases above ``the same'' manifold?
  \item How many smooth structures does a topological manifold carry in dimension $1,2,3$, above $4$, and equal to $4$?
  \item Why are tangent spaces not defined as planes in an ambient space?
\end{enumerate}
\textbf{Concepts.}
\begin{itemize}
  \item $C^\star$ compatibility, $C^k$, smooth, analytic and complex atlases (\cref{def:df-star,def:df-ck,def:df-cr}); the theorem on smooth sub-atlases (\cref{thm:df-whitney}).
  \item $C^k$ manifold (\cref{def:df-manifold}); incompatible atlases (\cref{rem:df-incompat,ex:df-line}).
  \item Differentiable maps between manifolds and their well-definedness (\cref{def:df-map,prop:df-welldef}); the manifold philosophy (\cref{int:df-philosophy}).
  \item Diffeomorphism (\cref{def:df-diffeo,ex:df-line-diffeo}); levels of structure (\cref{sec:df-levels}).
  \item Classification: dimensions $1,2,3$, above $4$, $4$ (\cref{thm:df-low,thm:df-high,thm:df-four}); intrinsic tangent spaces (\cref{int:df-embedding}).
\end{itemize}
\end{keyconcepts}

\flushside
\end{paracol}
